\section{Global Economy and International Organisations}

%%% From the Foreign Trade Module



\begin{frame}{The World is large---and all we have}
    \mycomment{The World Economy is the only `closed economy' we have. 
    
    There is production and consumption, and there is global growth---
    the world economy is not a zero sum game. 
    
    There are also fundamental balances: every buyer has a seller, every import is also an export.
    }
    \showmytikz{shared/tikz/tikz-globe-europe-africa}
\end{frame}




\begin{frame}{International Law}
    \mycomment{
    International law exists, but not in the same way as national, codified, positive law exists.

    There is natural law, which is considered to exist prior to set and codified law. Among this are the human rights, thought to exists independent of formulation. Set law that is agreed upon or issued by law makers is called positive law. 
    }
    \mytext{
        \begin{itemize}
        \item There is no World State
        \item International Law is contract law
        \item Concept of Natural Law vs. positive law
        \item Declaration of Human Rights
        \item No enforcement authorities
        \item International Law of the Seas
        \item UN System
        \item OECD
        \item Local legal systems: EU Law
    \end{itemize}
    }
\end{frame}

\begin{frame}{The Open Economy}
    \mycomment{
    No man is an island, goes the adage. The world itself is the only closed economy we have. Approaches to close off economies from the rest of the world have typically been desastrous. 

    Let us consider again the model \myfigref{fig:flowmodel} of the open economy with labour markets, markets for goods and services, and capital markets, all influenced by economic policy of the state.
    }
    \showmytikz[1,2]{shared/tikz/tikz-flowmodel-open}
\end{frame}

\begin{frame}{The Open Economy}
    \mycomment{
    Now our model economy comes into contact with another economy, and we may well assume this to be just one, that is the \textbf{world market.}
    }
    \showmytikz[1,2,3]{shared/tikz/tikz-flowmodel-open}
\end{frame}


\begin{frame}{The Open Economy}
    \mycomment{
    As the economy opens up to the world market, goods traded aboard ships enter, and goods traded inside are offered on world markets and shipped out where competitive. So the world market and the local market combine, if there are no trade barriers. The same is true for capital markets unless there are capital controls.

    Slightly different is the situation on the labour markets, as movement of people across borders is more difficult. 
    
    However, as we know from Walras' Law, if two of three markets are in equilibrium, so must be the third.    
    It follows that when there is no capital control and free trade, we expect labour markets to converge as well. 
    }
    \showmytikz[1,2,3,4]{shared/tikz/tikz-flowmodel-open}
\end{frame}




\begin{frame}{International Law}
    \mycomment{
    International (positive) law originates from agreements between states. There is no super national authority that can force states under its will. All international cooperation is thus voluntary. However, there are standards that are agreed upon, constituting some kind of international common, or positive, law.
    }
    \mytext{
    \begin{itemize}
        \item Hugo Grotius 1583-1645
        \item States as subjects, not firms or people
        \item People to minor extent (per Art. 1 UN Charta)
        \item Also organisations of states
        \item UN is such an organisation
        \item UN charter
        \item UN declaration of human rights (non binding!)
        \item conventions and treaties (Geneva Convention Relating to the Status of Refugees)
        \item Specialised Agencies: IMF, World Bank, ILO etc.
    \end{itemize}    
    }
\end{frame}


\begin{frame}{UN System}
    \showmytree[proper]{../shared/tikz/forest-unsystem}
    \mycomment{
         The UN is an organisation with states as members. So it does not consitute another level of government (no \enquote{world state}).

         Groups of countries like the G7 (G8, G20) are just conferences of the respective heads of states. 
    }
\end{frame}

\begin{frame}{UN System}
    \showmytree[special]{../shared/tikz/forest-unsystem}
    \mycomment{
    IMF and World bank were actually founded 1944, a year before the UN, at Bretton Woods (\enquote{Bretton Woods institutions}).
    
    The International Monetary Funds provides funds (as the name suggests) to stabilise exchange rates. It provides financial assistance (lending) at lower-than-market interest rates to countries with problems with balance-of-payments difficulties.

    The World Bank provides credits to large development programmes. Thus other than the name suggests, it is \textit{not} a \enquote{world central bank}. Its lending arm is the IBRD (International Bank for Reconstruction and Development). 

    Both institutions were founded as part of the Bretton Woods System of fixed exchange rates after the Second World War. This system of fixed exchange rates was abandoned in 1973. So today the role of the IMF is not to help guarantee fixed, but to help stabilise floating exchange rates. 
    }
\end{frame}

\begin{frame}{UN System}
    \mycomment{
    There are a number of \enquote{other entities and bodies} of the United Nations, notably the United Nations University, a think tank, headquartered in Tokyo.
    }
    \showmytree[other]{../shared/tikz/forest-unsystem}
\end{frame}

\begin{frame}{UN System}
    \mycomment{
    A number of organisations are \enquote{related} to the United Nations, but not all UN members are automatically members. The most important of these is probably the WTO, the World Trade Organisation, headquartered in Geneva.
    }
    \showmytree[related]{../shared/tikz/forest-unsystem}
\end{frame}

\begin{frame}{EU System}
    \showmytree{../shared/tikz/forest-eusystem}
    \mycomment{
    The constitutional text are the Treaty of the European Union (TEU) and the Treaty of the Functioning of the European Union (TFEU).

    There are three types of EU core tools to make and guide policy (legislation types):

    \begin{itemize}
        \item Regulations: directly binding
        \item Directives: must be implemented by member states
        \item Communications: non-binding
    \end{itemize}
    }
\end{frame}

\begin{frame}{EU System}
    \showmytree[1]{../shared/tikz/forest-eusystem}
    \mycomment{
        The four decision making organs are core institutions of the EU. 
        
        The Commission has the right of initiative to propose new laws, which are passed by the Parliament and the Council of the Europen Union jointly. 
        
        The Commission also acts as executive. The European Council defines the general political direction. 
    }
\end{frame}

 \begin{frame}{EU System}
    \showmytree[2]{../shared/tikz/forest-eusystem}
    \mycomment{
        Among the other organs of the EU are the Court of Justice with two courts, the European Central Bank, and the European Court of Auditors. 

        The European Central Bank is an independent central bank setting the monetary policy for the EU.

    }
\end{frame}

  \begin{frame}{EU System}
    \showmytree[3]{../shared/tikz/forest-eusystem}
    \mycomment{
        Agencies of the European Union are in charge to coordinate many other fields of policy of the EU.
    }
\end{frame}

 

\begin{frame}{European Union - Competencies}
    \showmytikz[1]{../shared/tikz/tikz-eucompetencies}
    \mycomment{The competencies of the European Union are defined in Art. 2 TFEU, while Art. 3 lists the exclusive competencies explicitly.

    Competition rules are an exclusive competency as long as the single market is concerned. 
    
    Monetary policy is an exclusive competency in regards to the Eurozone countries. 
    
    Trade is the common commercial policy. There are exceptions, in particular in areas where all competencies remain with member states (e.g. military equipment). 
    
    Fishery is an exclusive competency, but only as far as it is regulated by the common fishery policy. 
    }
\end{frame}

\begin{frame}{European Union - Competencies}
    \showmytikz[2,3]{../shared/tikz/tikz-eucompetencies}
    \mycomment{
    In other areas, competencies are either shared, or the EU may play a supportive role.
    }
\end{frame}


\begin{frame}{OECD System}
    \showmytree[1,2]{../shared/tikz/forest-oecdsystem}
    \mycomment{
     The Organisation for Economic Cooperation and Development is not an organisation of states, but a permanent conference and its corresponding organisation. It is headquartered in Paris. The member states are the industrialised countries (only), so often \enquote{OECD member states} is used when comparing industrialised countries only. These also report to the OECD, making the OECD a valuable source of country and policy information, particularly for comparisons.

     Note the International Energy Agency IEA is a body of the OECD, not of the UN.
    }
\end{frame}



\begin{frame}{GDP in USD, 2025}
    \mycomment{
    GDP of countries are the main measure for the economic importance of a country. However, it does not count stocks. There could be countries (Switzerland? UK? Saudi Arabia?) that hold many assets, but produce relatively little. The income generated (for example in the US) is counted where it is produced - independent of the fact that it may end up being owned abroad. 
    }
    \mytable{
        \begin{tabular}{rlrl}
            \toprule
            \def\myblock{fill=black}
            & & Trillion USD &  \\
            & World & 110.0 &  \\
           \midrule
            1st & USA  & 30.6   & \myblock[0.4ex]{30.6}\\
                & EU   & 21.1   & \myblock[0.4ex]{21.1}\\
            2nd & China & 19.4  & \myblock[0.4ex]{19.4}\\
            3rd & Germany & 5.0 & \myblock[0.4ex]{5.0}\\
            4th & Japan & 4.3   & \myblock[0.4ex]{4.3}\\
            5th & India & 4.1   & \myblock[0.4ex]{4.1}\\
            $\vdots$ & $\vdots$ & $\vdots$ & \\
            \bottomrule
        
        \end{tabular}

        Note: 1 Trillion $= 10^{12}$
    
        Source: \cite{IMFdatamapper}
    
    }
\end{frame}

\begin{frame}{GDP Growth Rates, 2026}
    \mycomment{
    Growth rates are interesting, because they give a hint how the absolute ranking may change. Refer to the `rule of 72'. At a growth rate of 4\%, any value will approximately double after $72/4=18$ periods.

    Compare Table \ref{tab:yields}
    }
    \mytable{\label{tab:gdpgrowhrates}
    \begin{tabular}{rlrl}
        \toprule
         &   &        &   \\
       GDP & World & 3.0      & \myblock[2ex]{3.0} \\
       \midrule
        1st & USA  & 2.1     & \myblock[2ex]{2.1}\\
            & EU   & 1.4     & \myblock[2ex]{1.4}\\
        2nd & China & 4.2    & \myblock[2ex]{4.2}\\
        3rd & Germany & 0.9  & \myblock[2ex]{0.9}\\
        4th & Japan & 0.6    & \myblock[2ex]{0.6}\\
        5th & India & 6.2    & \myblock[2ex]{6.2}\\
        $\vdots$ & $\vdots$ & $\vdots$ &  \\
        \bottomrule
    \end{tabular}

    Source: \cite{IMFdatamapper}
    
    }
\end{frame}


\begin{frame}{GDP in USD per Capita 2025}
    \mycomment{The absolute GDP does not tell us about the standard of living or grade of development of a country. There are countless measures, but the easiest to handle is GDP per capita.
    
    This does of course not take into account how evenly or unevenly this total income may be distributed. There may be much inequality, and much poverty, in the `richest' of nations.
    }
    \mytable{
        \begin{tabular}{lrl}
            \toprule
             & TUSD/a & \\
            \midrule
            Switzerland & 110       & \myblock[0.18ex]{110}\\
            USA &  89               & \myblock[0.18ex]{89}\\
            Germany & 59            & \myblock[0.18ex]{59}\\
            Europe & 42             & \myblock[0.18ex]{42}\\
            Japan & 34              & \myblock[0.18ex]{34}\\
            Poland & 28             & \myblock[0.18ex]{28}\\
            \textbf{World} & 15     & \myblock[0.18ex]{15}\\
            Mexico & 14             & \myblock[0.18ex]{14}\\
            China & 14              & \myblock[0.18ex]{14}\\
            India & 3               & \myblock[0.18ex]{3}\\
            \bottomrule
        \end{tabular}
        
         Source: \cite{IMFdatamapper}
    }
\end{frame}

\begin{frame}{Government Debt, 2026}
    \mycomment{
    Gross general government debt in percent of GDP. High debt is not a problem per se (Japan), but it makes a country vulnerable to rising interest rates: the relevant comparison is the interest rate $r$ vs.\ the growth rate $g$ (see previous slide).
    }
    \mytable{
    \begin{tabular}{rlrl}
        \toprule
         &   & \% GDP &   \\
       GDP & World & $\approx$95 & \myblock[2ex]{1.9} \\
       \midrule
        1st & USA       & 125.8  & \myblock[2ex]{2.52}\\
            & Euro area &  87.8  & \myblock[2ex]{1.76}\\
        2nd & China     & 106.9  & \myblock[2ex]{2.14}\\
        3rd & Germany   &  64.6  & \myblock[2ex]{1.29}\\
        4th & Japan     & 204.4  & \myblock[2ex]{4.09}\\
        5th & India     &  83.4  & \myblock[2ex]{1.67}\\
        $\vdots$ & $\vdots$ & $\vdots$ &  \\
        \bottomrule
    \end{tabular}

    Source: \cite{IMFfiscalmonitor} 
    
    }
\end{frame}

\begin{frame}{Real Government Bond Yields, 10 Years, Oct 2026}
    \mycomment{
    Real yield = nominal 10-year yield minus current inflation (ex post). This is the real cost of new government borrowing. If the real yield is below real GDP growth ($r<g$), the debt ratio can fall even with small primary deficits.

    Compare Table \ref{tab:gdpgrowhrates}
    }
    \mytable{\label{tab:yields}
    \begin{tabular}{rlrrl}
        \toprule
         &   & nominal & real &   \\
       GDP & World &  --  &  --  &  \\
       \midrule
        1st & USA       & 4.9 & 1.2 & \myblock[2ex]{1.2}\\
            & Euro area &  -- &  -- & \\
        2nd & China     & 1.7 & 0.8 & \myblock[2ex]{0.8}\\
        3rd & Germany   & 3.3 & 0.5 & \myblock[2ex]{0.5}\\
        4th & Japan     & 2.8 & 0.7 & \myblock[2ex]{0.7}\\
        5th & India     & 7.4 & 1.9 & \myblock[2ex]{1.9}\\
        $\vdots$ & $\vdots$ & $\vdots$ & $\vdots$ & \\
        \bottomrule
    \end{tabular}

    Sources: \cite{BloombergRates} \cite{TEinflation} \cite{ChinaBond};
    CPI latest available (Sep/Oct 2026)
    
    }
\end{frame}

\begin{frame}{Balance of Payments}
    \mytext{
    The Balance of Payments (BoP) is an account that records all transactions crossing the border of an economy.

    It corresponds to accounts of the System of National Accounts (SNA), where positions appear in \enquote{rest of the world}-accounts:
    \begin{itemize}
        \item from current accounts: external account of goods and services, external accounts of income, etc.: \textbf{current account of BoP}
        \item from accumulation accounts: external financial and capital accounts        
    \end{itemize}
    }
    \mycomment{
    The BoP is thus an account that \enquote{collects} all data from the SNA that involves trans-border transactions. Again, there is a credit and a debit side, depending on the sign of the net position.

    Due to double book-keeping, the BoP is also always in balance. Anything crossing the border is double-booked.
    }
\end{frame}

\begin{frame}{Balance of Payments}
    \showmytikz{../shared/tikz/tikz-boptikz}
    \mycomment{

    Depending on whether the Trade Balance is positive or negative, it will appear on the left or right of the BoP. The same applies for the other Accounts.

    \textbf{The Balance of Payments is always level.} There can be deficits in the Current Account, or in the Financial and Capital Accounts. But because every transaction is recoreded twice, the balance is always zero. 

    Under fixed exchange rates, the central bank must intervene to hold the peg; these interventions appear in the reserve account. Under floating rates, reserve changes are small and the current account is matched mainly by private capital flows.
    

    Contrary to what the names suggest, the Balance of Payments is not a balance, but shows \textbf{flow quantities: inflows, outflows over a period.}
    }
    
\end{frame}

\begin{frame}{Imports and Exports in the BoP}
    \showmytikz{../shared/tikz/tikz-bop-entries}
    \mycomment{
    Export of goods and services is either paid for in cash ($\Delta R$), or made possible through an import of a liability. 
    
    Thus normally, \textbf{exports of goods} correspond to \textbf{export of capital.} (Better speak of \enquote{net acquisition of financial assets} or \enquote{net incurrence of liabilities}).

    Under the BPM6/7 sign convention, the financial account shows net lending: a current account surplus is matched by a financial account balance of the same sign.

    A deficit in the current account can thus correspond to a deficit in the capital account, as excessive import of goods is accompanied by the export of a liability, which is an import of capital.

    Foreign capital inflow correspond to current account deficits; which causes what is another question altogether not settled by the identity. 

    Note that currency reserves also include foreign government securities, in particular those of the US (USTs).
    }

\end{frame}

\begin{frame}{Balance of Payments in Balance}

    \mycomment{
        There is sometimes confusion: a trade deficit or surplus means an imbalance in the current account, or the trade balance as part of the current account.

        It follows that imports not covered by exports ($Im>Ex$) can only appear when capital import happens ($K_{Im}>K_{Ex}$); that is the credit the exporter gives the importer.

        The Reserve Account consists of cash, but also of e.g. government securities. This is often UST (T-Bills), in particular in USD.

        The rest of the world is financing the US trade deficit through capital export into the US, often by buying highly liquid USTs. Weaker currencies, states or governments do not hold that privilege.
    }
        
    \mytext{
        The BoP is \textbf{always, by definition, in balance.} 
        
        \begin{equation}
            \colorbox{yellow!40}{$Ex-Im-K_{Ex}+K_{Im}-\Delta R = 0$}
        \end{equation}

        To see this, look at the sub-balances:
        
        \begin{equation}
             \text{\scriptsize (Current A.)} - \text{\scriptsize (Fin \& Cap. A.)} =  \text{\scriptsize (Reserve A.)  }
        \end{equation}


    }
\end{frame}


\begin{frame}{Balance of Payments (IMF)}
    \mycomment{
    This tree represents the Balance of Payments (account) as defined by the IMF \cite{imfbpm7}. 

    All these \enquote{balances} register \textit{flow values} -- other than their name suggests. 

    The BoP consists of two accounts: 
    \begin{itemize}
        \item the \textbf{current account}, where current transactions across borders are recorded; 
        \item \textbf{capital and financial account} is an accumulation account. The \textbf{reserve account} is a balance, and it records a surplus or deficit in currencies.
    \end{itemize}
    }
    \showmytree[1,3]{../shared/tikz/forest-bop-imf.tex}
\end{frame}

\begin{frame}{BoP: Current Account (IMF)}
    \mycomment{
    The Current Account consists of the \textbf{trade balance}, which is a positive or negative balance of the account for goods and services.

    In addition, the Current Account also registers cross-border payments for labour (compensation of employees) and interest payment (investment income). 

    And lastly, it also registers current transfers, among which transfers as development aid, etc.
    }
    \showmytree[only 1,2]{../shared/tikz/forest-bop-imf.tex}
\end{frame}

\begin{frame}{BoP: Capital and Financial Account  (IMF)}
    \mycomment{
    The Capital and Financial Account records changes of ownership and hence financial transfers, most importantly cross-border investment: foreign direct investment (FDI) and portfolio investment.

    This account also comprises the Reserve Account, which is an account at the Central Bank registering the \textit{change} of reserve assets.  
    }
    \showmytree[only 3,4]{../shared/tikz/forest-bop-imf.tex}
\end{frame}

\begin{frame}{Trade as Part of Global GDP}
    \mycomment{Trade is a large part of global GDP. That is, a large part of what is produced world wide is traded. 
    
    That does of course not measure precisely the contribution of trade to GDP (and GDP growth), which is hard to measure.
    }
    \mytext{
    \begin{itemize}
        \item Global trade 2024: ca. 33 Trillion USD
        \item 8-9 trillion USD services
        \item 24 trillion USD goods
    \end{itemize}
    Source: \href{https://unctadstat.unctad.org/datacentre/dataviewer/US.TradeMerchTotal}{UNCTAD}

    As global GDP stands at 110 Trillion USD, trade is about \textbf{30\%.}
    }
\end{frame}

\begin{frame}{Trade Indicators}
\mycomment{
    The (trade) openness of a country is measured by setting the sum of exports and imports relative to gross domestic product.
    
    It does not, however, tell us what type of goods are imported and exported: raw materials, or consumer goods, or energy? And are these all imported or exported to a plethora of partners, or are there dependencies on singular trading partners?
    
    One can of course also put just imports or just exports in relation to GDP. This gives you the Import or Export Penetration Rate.

    The Dependency Rate measures how much of a good $X$ (raw material, energy) is imported relative to how much of it is consumed locally. 

    }
\mytext{
    Trade Openness:
    \begin{equation}
        TO = \frac{\left|{Ex}\right|+\left|{Im}\right|}{{Y}}
    \end{equation}  

    Import Ratio:
    \begin{equation}
        IPR = \frac{Im}{Y}
    \end{equation}

    Import Penetration Ratio:
    \begin{equation}
        IPR= \frac{Im}{Y-Ex+Im}
    \end{equation}

    Dependency Rate:
    \begin{equation}
        Dr_X=\frac{Im_X}{C_X}
    \end{equation}
}
\end{frame}

\begin{frame}{Trade Indicators}
\mycomment{
Openness varies widely from country to country. But larger countries are of course `less open' than small countries. This does not say anything about their trade policy, but only about the part of GDP that is traded. 

Remember that trade openness counts imports and exports, thus every cross-border trade is counted twice. To measure the influence on GDP, one has to look at the trade balance in accordance with the fundamental equation of macroeconomics (eq.~\ref{eq:fund}).

}
\mytable{\scriptsize
    \begin{tabular}{lrl}
        \toprule
         Brazil & 38.8               & \myblock[0.18ex]{38.8}\\
         Chile & 75.2                & \myblock[0.18ex]{75.2}\\
         China & 37.4                & \myblock[0.18ex]{37.4}\\
         Congo Rep. & 98.2           & \myblock[0.18ex]{98.2}\\
         EU & 103.7                  & \myblock[0.18ex]{103.7}\\
         Germany & 100              & \myblock[0.18ex]{100}\\
         High Income & 67.4          & \myblock[0.18ex]{67.4}\\
         Hong Kong & \textbf{352.9}          & \\
         Indonesia & 45.5            & \myblock[0.18ex]{45.5}\\
         Least developed (UN) & 51.2 & \myblock[0.18ex]{51.2}\\
         Mexico & 88.8               & \myblock[0.18ex]{88.8}\\
         UK & 69.2                   & \myblock[0.18ex]{69.2}\\
         USA & 27.3                  & \myblock[0.18ex]{27.3}\\
         World & 62.4                & \myblock[0.18ex]{62.4}\\
         \bottomrule
        \bottomrule
    \end{tabular}
    
    Data for 2022 \cite{worldbankdatabank}
}
\end{frame}

\begin{frame}{Imports Ratio to GDP}
    \mycomment{
    As presented by Kristalina Georgieva, Managing Director, IMF on April 17th, 2025 \cite{IMFgeorgieva2025}. Large economies are typically much less open than smaller ones. 
    
    Protectionism affects smaller countries more than larger ones.
    }
    \showmybitmap[https://www.imf.org/en/news/articles/2025/04/17/sp041725-spring-meetings-2025-curtain-raiser]{../shared/bitmaps/6-curtain-raiser-charts-sms-final.jpg}
\end{frame}

\begin{frame}{WTO and GATT}
    \mycomment{
    In the economic realm, many frameworks exist to make international trade more safe and secure. 
    }
    \mytext{
    The General Agreement on Tariffs and Trade (GATT) is from 1949. 
    The WTO was founded 1995. China joined 2001.
     \begin{itemize}
        \item WTO $\neq$ UN, fewer members
        \item WTO $\leftarrow$ GATT
        \item GATT still in existence
        \item Dispute Settlement Body
        \item Appellate Body: paralysed since 2019
        \item Art.\ 1 GATT: Most Favoured Nation (MFN) Principle
        \item Exception: customs union, FTAs
    \end{itemize}   
    }
\end{frame}

\begin{frame}{Most Favoured Nation Clause}
    \mycomment{
    The MFN clause: charge every trading partner the lowest tariff you charge to anyone else.
    }
    \mytext{
    \begin{quote}
        {\scriptsize With respect to customs duties and charges of any kind imposed on or in connection with importation or exportation or imposed on the international transfer of payments for imports or exports, and with respect to the method of levying such duties and charges, and with respect to all rules and formalities in connection with importation and exportation, and with respect to all matters referred to in paragraphs 2 and 4 of Article III,
        any advantage, favour, privilege or immunity granted by any contracting party to any product originating in or destined for any other country shall be accorded immediately and unconditionally to the like product originating in or destined for the territories of all other contracting parties.} \cite{gatt1947}
    \end{quote}
    }
\end{frame}


\begin{frame}{WTO and Trade Blocs}
    \mytext{
    GATT/GATS/TRIPS $\rightarrow$ WTO, Geneva: 
    \begin{itemize}
        \item `most favoured nation' clause
        \item ratchet principle (tariffs can only go down)
        \item differentiation only justifiable by blocs
        \item result: surge in blocs
    \end{itemize}
    }
\end{frame}


\begin{frame}{WTO and Free Trade Areas}
    \showmytikz[1]{../shared/tikz/tikz-wto}
    \mycomment{
        The WTO is an organisation of member states based in Geneva. It runs, among others, a dispute settlement body. It came to replace the (still valid!) GATT, the general agreement on tariffs and trade.

        Within the WTO, the MFN most favoured nation principle holds: you are not allowed to charge anyone a tariff higher than you charge to the most favoured nation (there are exceptions).

        Exceptions are only possible if you are within a free trade area. 
    }
\end{frame}

\begin{frame}{WTO and Free Trade Areas}
    \showmytikz[1,2]{../shared/tikz/tikz-wto}
    \mycomment{
        Since you can only positively discriminate if you are in a free trade area, free trade areas sprang up.

        The degree of economic integration within these areas varies widely. The simplest is just a free trade area, so no tariffs apply. But it can reach all the way to an integration like the EU, with a common currency, common standards, and common legal bodies.     

        (Picture is illustrative: USMCA has exactly three members, USA, Mexico, and Canada. Mercosur and EU have more members than shown, etc.)
    }
\end{frame}


\begin{frame}{Trade Blocs / PTAs }
    \mycomment{
        Preferential trade agreements (PTAs) can be more or less intense. In a free trade agreement, you still need to file customs documents. In a customs union, you may still encounter different regulations. In a common market, you may still encounter different economic policies. In an economic union, you may still have to deal with exchange rates. So an economic union must include a common market, and so forth.
    
        While every currency union is also an economic union, is a common market, etc., geographically it works just the other way around: a currency union may be embedded in a (larger) economic union, which is within a (larger) common market, etc. 

        In reality, these circles do not nest perfectly, but rather overlap (at the edges).
    }
    \showmytikz[2]{../shared/tikz/tikz-ftas-venn}
\end{frame}

{%
    \beameronly{
        \setbeamertemplate{background}{%
                \drawmap[tikz={remember picture, overlay}, at=current page.center,context,
                       width=.9\paperwidth, height=.8\paperheight,
                       style={fill=yellow!20, draw=blue!40}]{%
                Benin,Cabo Verde,Côte d’Ivoire,Ghana,Guinea,Guinea Bissau,Liberia,Nigeria,Senegal,Sierra Leone,The Gambia,Togo}
        }
    }
\begin{frame}{Free Trade Zones Examples}
    \mycomment{
    There are plenty of free trade areas and monetary cooperation areas. Sometimes they go by French or English abbreviations. Negotiations typically take years, and scope, contracts, names and memberships may change.

    An overview and date can be found at \cite{wto_rtadb}.
    }
    \mytext{
    \begin{itemize}
        \item EMU, EU, EEA, EUCU
        \item Mercosur, CACM etc.
        \item AFTA (ASEAN)
        \item JEFTA (EU-Japan), CETA (EU-Canada)
        \item CEDEAO $=$ ECOWAS

    \end{itemize}
    }
\end{frame}
}



{%
    \beameronly{
        \setbeamertemplate{background}{%
                \drawmap[tikz={remember picture, overlay}, at=current page.center,context,
                       width=.9\paperwidth, height=.8\paperheight,
                       style={fill=yellow!20, draw=blue!40}]{Benin, Burkina Faso, Côte D'Ivoire, Guinea-Bissau, Mali, Niger, Senegal, Togo,
                       Cameroon, Central African Republic, Chad, Equatorial Guinea, Gabon, Republic of the Congo}
        }
    }
    \begin{frame}{Currency Blocs}
        \articleonly{
        \begin{figure}[ht!]
            \centering
            \drawmap[context,
                       width=0.61\textwidth, height=.8\paperheight,
                       style={fill=yellow!20,draw=black}]{Benin, Burkina Faso, Côte D'Ivoire, Guinea-Bissau, Mali, Niger, Senegal, Togo,
                       Cameroon, Central African Republic, Chad, Equatorial Guinea, Gabon, Republic of the Congo}
            \caption{UEMOA and CEMAC}
            \label{fig:uemoa_cemac}
        \end{figure}
        }
        \mytext{
        \begin{itemize}
            \item EMU (Euro)
            \item UEMOA=WAEMU and CEMAC (tied to \euro ) \myfigref{fig:uemoa_cemac}
        \end{itemize}
    
        Unilateral peg
        \begin{itemize}
            \item Kosovo (uses Euro outright)
            \item Zimbabwe (?)
            \item Morocco
            \item Argentina..?
        \end{itemize}
        }
    \end{frame}
}

{%
    \beameronly{
        \setbeamertemplate{background}{%
        \drawmap[tikz={remember picture, overlay}, at=current page.center,context,
                       width=.9\paperwidth, height=.8\paperheight,
                       style={fill=yellow!20, draw=blue!40}]{Argentina, Bolivia, Brazil, Paraguay, Uruguay}
        }
    }
    \begin{frame}{Regional, Economic and Political Cooperation}
        \articleonly{
            \begin{figure}[ht!]
                \centering
                \drawmap[context,
                       width=.61\textwidth, height=.8\paperheight,
                       style={fill=yellow!20,draw=black}]{Argentina, Bolivia, Brazil, Paraguay, Uruguay}   
                        \caption{Mercosur}
                \label{fig:mercosur}
            \end{figure}
        }
        \mytext{
         \begin{itemize}
            \item Working Groups (G7, G20)
            \item Local Working Groups (E3)
            \item Currency union(s), customs union(s)
            \item EU, EEA, ECU
            \item MERCOSUR \myfigref{fig:mercosur}
            \item ASEAN
            \item USMCA
            \item Economic Groups 
            \item OPEC
            \item BRICS
            \item ...
        \end{itemize}   
        }
        \mycomment{
        The intensity of the cooperation can and does vary widely. Most of them are mere conferences, such as the G7. They may or may not have a secretariat. %Some , that may not be more than an office in New York at the UN. 
    
        In others, most prominently the European Union, member states transfer parts of their sovereignty to multi-national bodies with large institutions---and budgets.
        }
    \end{frame}
}

{%
    \beameronly{
        \setbeamertemplate{background}{%
           \drawmap[tikz={remember picture, overlay}, at=current page.center,context,
           width=.9\paperwidth, height=.8\paperheight,
           style={fill=yellow!20, draw=blue!40}]{
                       Brunei Darussalam,Cambodia,Indonesia,Lao PDR ,Malaysia,Myanmar,Philippines,Singapore,Thailand,Vietnam 
                       }
            }
    }
    \begin{frame}{Asia/Pacific: ASEAN}
        \articleonly{
        \begin{figure}[ht!]
            \centering
           \drawmap[context,
           width=.61\textwidth, height=.8\paperheight,
           style={fill=yellow!20,draw=black}]{
                       Brunei Darussalam,Cambodia,Indonesia,Lao PDR ,Malaysia,Myanmar,Philippines,Singapore,Thailand,Vietnam 
                       }            
            \caption{ASEAN}
            \label{fig:asean}
        \end{figure}
        }
        \mycomment{
        Asian countries invoice heavily in US Dollar, even when trading with China. Economic integration in the area is comparatively weak. 
        
        Japan in particular is a single country with a single currency (Japanese Yen, which means \enquote{circle}) without a \enquote{Yen zone}. 
        
        The Chinese economy is huge, but capital flows are typically restricted; the currency is the \enquote{Renminbi}, which means \enquote{currency of the people}, with the unit Yuan, which means \enquote{unit}. The cognate with Yen is no coincidence.
        }
        \mytext{
        ASEAN Association of Southeast Asian Nations
        \begin{itemize}
            \item AFTA Free trade area
            \item AIA Asean Investment Area
        \end{itemize}
    
        }
    \end{frame}
}

{%
    \beameronly{
        \setbeamertemplate{background}{%
           \drawmap[tikz={remember picture, overlay}, at=current page.center,context,
           width=.9\paperwidth, height=.8\paperheight,
           style={fill=yellow!20, draw=blue!40}]{
                       Botswana,Eswatini,Lesotho,Namibia,South Africa% 
                       }
            }
    }
    \begin{frame}{Africa: AEC, regional blocs}
        \articleonly{
        \begin{figure}
            \centering
           \drawmap[context,
           width=.61\textwidth, height=.8\paperheight,
           style={fill=yellow!20,draw=black}]{
                       Botswana,Eswatini,Lesotho,Namibia,South Africa% 
                       }            
            \caption{AEC}
            \label{fig:aec}
        \end{figure}
        }
        \mycomment{
        Africa with 54 states has a plethora of international organisations. Much of francophone Africa is closely tied to the Euro area, while chiefly anglophone South Africa (the Rand zone) and East Africa lean more towards the US Dollar.
        }
        \mytext{
            \begin{itemize}
            \item AEC African Economic Community \myfigref{fig:aec}
            \begin{itemize}
                \item ECOWAS/CEDEAO Economic Community of West African States
                \item WAEMU/UEMOA West African Economic and Monetary Union (Franc CFA)
                \item ECCAS/CEEAC Economic Community of Central African States    
                \item EMCCA/CEMAC Economic and Monetary Community of Central Africa
                \item ...
                \item SADC Southern African Development Community
                \item SACU Southern African Customs Union
                \item ...
            \end{itemize}
        \end{itemize}
        }
    \end{frame}
}
